\subsection{Gaussian filters under finite coordinate changes}
\label{sec:gaussian_reindex}

% Source: polynomial-PEPS manuscript (September 24, 2026), lem:reset,
% pinned peps-02-information-adc7f124.tex, lines 426--452, at
% openai/math@adc7f1241b42e322a6451854ab7e4b4c146bf78a.
% This is the coordinate-transport interface for the Gaussian integral.
% Regional threeSplit/sheetSwap identifications are downstream in TNLean.
% No locality estimate or complete reset theorem is asserted here.

Let $n,m$ be finite sets and let $e:n\to m$ be a bijection. For a
complex matrix $A$ indexed by $n$, define its reindexing $R_e(A)$ by
\begin{align}
    R_e(A)_{ab}&=A_{e^{-1}(a),e^{-1}(b)}\qquad(a,b\in m).
    \label{eq:gaussian_reindex_coordinates}
\end{align}
Thus the same inverse bijection acts on both indices, with the entries
unchanged. The map $R_e$ is an algebra isomorphism and a continuous
complex-linear equivalence. The finite sets may be empty.
For any complex matrix $H$, write $U_H(t)=e^{itH}$; Hermiticity is not
assumed in the coordinate identities below.

\begin{theorem}[Exponential covariance]
    \label{thm:gaussian_reindex_exp}
    \lean{Matrix.reindex_exp, Matrix.reindex_hermitianUnitaryPath}
    \leanok
    For all complex matrices $A,H$ indexed by $n$ and all $t\in\R$,
    \begin{align}
        R_e(e^A)&=e^{R_e(A)}.
        \label{eq:gaussian_reindex_exp}
    \end{align}
    \begin{align}
        R_e(U_H(t))&=U_{R_e(H)}(t).
        \label{eq:gaussian_reindex_path}
    \end{align}
\end{theorem}
\begin{proof}\leanok
    The continuous algebra isomorphism $R_e$ preserves the exponential
    series. Apply this identity to $itH$ and use complex linearity to
    obtain~(\ref{eq:gaussian_reindex_path}).
\end{proof}

\begin{theorem}[Covariance of the matrix integral]
    \label{thm:gaussian_reindex_integral}
    \lean{Matrix.reindex_integral}
    \leanok
    Let $\mu$ be a measure on a measurable space $X$ and let
    $F:X\to\C^{n\times n}$ be any function. Then
    \begin{align}
        R_e\!\left(\int_X F(x)\,d\mu(x)\right)
            &=\int_X R_e(F(x))\,d\mu(x).
        \label{eq:gaussian_reindex_integral}
    \end{align}
    The integrals use the convention that a nonintegrable function
    has integral zero. No integrability assumption is needed.
\end{theorem}
\begin{proof}\leanok
    A continuous linear map commutes with the integral of an integrable
    function. Both $R_e$ and its inverse are continuous linear maps, so
    $F$ is integrable exactly when $R_e\circ F$ is integrable. In the
    remaining case both integrals are zero, and $R_e(0)=0$.
\end{proof}

\begin{theorem}[Covariance of the two-generator integrand]
    \label{thm:gaussian_reindex_integrand}
    \lean{GaussianFilter.reindex_intertwinerIntegrand}
    \leanok
    For all complex matrices $H',H,W$ indexed by $n$ and every real $t$,
    \begin{align}
        R_e\bigl(U_{H'}(t)WU_H(-t)\bigr)
            &=U_{R_e(H')}(t)R_e(W)U_{R_e(H)}(-t).
        \label{eq:gaussian_reindex_integrand}
    \end{align}
\end{theorem}
\begin{proof}\leanok
    \uses{thm:gaussian_reindex_exp}
    Apply multiplicativity of $R_e$ to the two products, then use
    (\ref{eq:gaussian_reindex_path}) for each exponential factor.
    The left generator remains $H'$ at time $t$ and the right generator
    remains $H$ at time $-t$.
\end{proof}

\begin{theorem}[Covariance of full and truncated Gaussian filters]
    \label{thm:gaussian_reindex_filters}
    \lean{GaussianFilter.reindex_gaussianIntertwiner,
        GaussianFilter.reindex_gaussianIntertwinerTruncated}
    \leanok
    \uses{def:gaussian_filter_integrals, def:gaussian_time_density}
    Let $h\ge0$, $T\in\R$, and let $H',H,W$ be arbitrary complex
    matrices indexed by $n$. For the centered Gaussian measure $\mu_h$
    of variance $h$, the actual full integral satisfies
    \begin{align}
        &R_e\!\left(\int_{\R}U_{H'}(t)WU_H(-t)\,d\mu_h(t)\right)
            \label{eq:gaussian_reindex_full_integral}\\
        &\qquad=\int_{\R}U_{R_e(H')}(t)R_e(W)U_{R_e(H)}(-t)\,d\mu_h(t).
            \notag
    \end{align}
    Its unrenormalized truncation satisfies
    \begin{align}
        &R_e\!\left(\int_{[-T,T]}U_{H'}(t)WU_H(-t)\,d\mu_h(t)\right)
            \label{eq:gaussian_reindex_truncated_integral}\\
        &\qquad=\int_{[-T,T]}U_{R_e(H')}(t)R_e(W)U_{R_e(H)}(-t)\,d\mu_h(t).
            \notag
    \end{align}
    Equivalently, $R_e(M_h(H',H;W))=M_h(R_e(H'),R_e(H);R_e(W))$,
    with the same identity for $M_h^T$. The variance, cutoff, and
    restricted Gaussian mass are preserved.

    These identities include $h=0$, when $\mu_0=\delta_0$ and the full
    integral is $W$. They also include $T<0$, when $[-T,T]$ is empty
    and both truncated integrals vanish. No Hermiticity, positive
    variance, nonnegative cutoff, or nonempty-index assumption is
    required.
\end{theorem}
\begin{proof}\leanok
    \uses{thm:gaussian_reindex_integral, thm:gaussian_reindex_integrand}
    Use Theorem~\ref{thm:gaussian_reindex_integral} first with $\mu_h$
    and then with its restriction to $[-T,T]$. In each case substitute
    (\ref{eq:gaussian_reindex_integrand}) inside the integral. The
    measure is unchanged, so there is no normalization factor.
\end{proof}
